% Einheit 1 von 5 – Rechteck, Quadrat, Umfang (bank/flaechen/e1.jsonl, zusammenbau v0.1)
%% TODO Zweigzeile Teil 1: was hier gelernt wird (Satz in Schülersprache aus dem Einheitstitel) – Einheit 1
\einheitenkopf[e1]{Einheit 1 von 5 · Rechteck, Quadrat, Umfang}
\zweigzeile{ab Kl. 5 · P10 oft}
\setcounter{aufgabe}{9}

%% TODO Titel: Ich-kann-Satz fehlt in der Bank (Platzhalter: Kettenname) – Nr. 10
%% TODO Anweisung: ein Satz über der ersten Teilaufgabe fehlt in der Bank – Nr. 10
\begin{aufgabe}{Gegeben – gesucht}
\begin{teile}
\teil Ein Rechteck hat die Fläche $A = 32$ m² und die Breite $b = 4$ m. Was ist gesucht? Kreise es in $A = a \cdot b$ ein.
\end{teile}
\end{aufgabe}

%% TODO Titel: Ich-kann-Satz fehlt in der Bank (Platzhalter: Kettenname) – Nr. 11
%% TODO Anweisung: ein Satz über der ersten Teilaufgabe fehlt in der Bank – Nr. 11
\begin{aufgabe}{Rechteck}
\begin{teile}
\teil Ein Garten bekommt einen Zaun. Fläche oder Umfang? \\ \kreuz{Fläche} \\ \kreuz{Umfang}
\end{teile}
%% TODO Darstellung für schwach fehlt in der Bank (flaechen-e1-k2-s1-v1; Abschnitt „Für schwache Schüler“)
\swz{Rechteck, $a = 8$ cm, $b = 3$ cm – Fläche und Umfang?}{}
%% TODO Darstellung für schwach fehlt in der Bank (flaechen-e1-k2-s1-v2; Abschnitt „Für schwache Schüler“)
\swz{Ein Beet ist $9$ m lang und $4$ m breit – Fläche und Umfang?}{}
%% TODO Darstellung für schwach fehlt in der Bank (flaechen-e1-k2-s1-v3; Abschnitt „Für schwache Schüler“)
\swz{Eine Tischplatte ist $7$ dm lang und $5$ dm breit – Fläche und Umfang?}{}
%% TODO Darstellung für schwach fehlt in der Bank (flaechen-e1-k2-s1-v4; Abschnitt „Für schwache Schüler“)
\swz{Ein Handybildschirm ist $11$ cm hoch und $6$ cm breit – Fläche und Umfang?}{}
%% TODO Darstellung für schwach fehlt in der Bank (flaechen-e1-k2-s1-v5; Abschnitt „Für schwache Schüler“)
\swz{Ein Kleinspielfeld ist $20$ m lang und $15$ m breit – Fläche und Umfang?}{}
\swfrage{Was bleibt gleich, was ändert sich?}
%% TODO Darstellung für schwach fehlt in der Bank (flaechen-e1-k2-s2-v1; Abschnitt „Für schwache Schüler“)
\swz{Rechteck, $a = 4{,}5$ cm, $b = 2$ cm – Fläche und Umfang?}{}
%% TODO Darstellung für schwach fehlt in der Bank (flaechen-e1-k2-s3-v1; Abschnitt „Für schwache Schüler“)
\swz{Quadrat, $a = 6$ cm – Fläche und Umfang?}{}
%% TODO Darstellung für schwach fehlt in der Bank (flaechen-e1-k2-s4-v1; Abschnitt „Für schwache Schüler“)
\swz{Ein Fünfeck hat die Seiten $3$ cm, $4$ cm, $5$ cm, $2$ cm und $6$ cm – Umfang?}{}
%% TODO Darstellung für schwach fehlt in der Bank (flaechen-e1-k2-s5-v1; Abschnitt „Für schwache Schüler“)
\swz{Rechteck, $A = 56$ cm², $a = 8$ cm – wie lang ist $b$?}{}
%% TODO Darstellung für schwach fehlt in der Bank (flaechen-e1-k2-s6-v1; Abschnitt „Für schwache Schüler“)
\swz{Rechteck, $u = 20$ cm, $a = 6$ cm – wie lang ist $b$?}{}
%% TODO Darstellung für schwach fehlt in der Bank (flaechen-e1-k2-s7-v1; Abschnitt „Für schwache Schüler“)
\swz{Quadrat, $A = 64$ cm² – Seitenlänge?}{}
%% TODO Darstellung für schwach fehlt in der Bank (flaechen-e1-k2-s8-v1; Abschnitt „Für schwache Schüler“)
\swz{Ein Rechteck ist doppelt so lang wie breit; die Breite ist $x$. Term für den Umfang?}{}
\swz[3]{Ein rechteckiger Bilderrahmen hat einen Umfang von $34$ cm, eine Seite ist $11$ cm lang. Wie lang ist die andere Seite? \hfill (P10 2018 OS)}{}
\end{aufgabe}

%% TODO Titel: Ich-kann-Satz fehlt in der Bank (Platzhalter: Kettenname) – Nr. 12
%% TODO Anweisung: ein Satz über der ersten Teilaufgabe fehlt in der Bank – Nr. 12
\begin{aufgabe}{aus Rechtecken zusammengesetzte Fläche (Summe)}
%% TODO Darstellung für schwach fehlt in der Bank (flaechen-e1-k3-s1-v1; Abschnitt „Für schwache Schüler“)
\swz{Ein L-förmiger Raum besteht aus zwei Rechtecken: $6$ m × $4$ m und $3$ m × $2$ m. Fläche?}{}
\end{aufgabe}

%% TODO Titel: Ich-kann-Satz fehlt in der Bank (Platzhalter: Kettenname) – Nr. 13
%% TODO Anweisung: ein Satz über der ersten Teilaufgabe fehlt in der Bank – Nr. 13
\begin{aufgabe}{Einheit wechseln vor dem Rechnen (cm und m gemischt)}
%% TODO Darstellung für schwach fehlt in der Bank (flaechen-e1-k4-s1-v1; Abschnitt „Für schwache Schüler“)
\swz{Rechteck, $2$ m lang, $45$ cm breit – Fläche in m²?}{}
\end{aufgabe}

%% TODO Titel: Ich-kann-Satz fehlt in der Bank (Platzhalter: Kettenname) – Nr. 14
%% TODO Anweisung: ein Satz über der ersten Teilaufgabe fehlt in der Bank – Nr. 14
\begin{aufgabe}{Rechteck – Fehler finden, Begründen, Anwendung, Darstellungswechsel}
\swz[3]{Ein Rechteck ist $6$ cm lang und $2{,}5$ cm breit. Tom rechnet den Umfang: \rechnung{u &= 6 \cdot 2{,}5 \\ u &= 15 \text{ cm}} Wo ist der Fehler?}{}
\swfrage{Warum rechnet man beim Rechteck Länge mal Breite? Erkläre mit Einheitsquadraten.}
\swz[3]{Ein Zimmer ist $4{,}2$ m lang und $3{,}5$ m breit. Laminat gibt es in Paketen zu $2{,}2$ m². Wie viele Pakete muss man kaufen?}{}
\swa{Zeichne ein Rechteck, dessen Fläche $A = 3 \cdot x$ ist, und beschrifte die Seiten.}{\rechenplatz[halb]{4}}
\end{aufgabe}

\uebersichtskasten{Fläche: Länge mal Breite. Umfang: einmal außen herum, alle Seiten addieren. \\ Rechteck a = 7 cm, b = 4 cm: A = 7 · 4 = 28 cm² \quad u = 2 · (7 + 4) = 22 cm \\ Quadrat a = 9 cm: A = 9 · 9 = 81 cm² \quad u = 4 · 9 = 36 cm \\ Rückwärts: Seite = A : andere Seite; Quadratseite = √A.}
